\section{From a balance response to a monetary decision}
\label{exec17:l2-monetary}

The deposit example exposes a difficulty that arises before any estimator is
chosen. A balance is a stock, its accumulation over time is an exposure, and
the financial benefit of that exposure depends on a valuation convention.
These quantities cannot be interchanged because they all contain a currency
unit. The assistant must obtain a complete statement of the business target:
does treasury want additional deposits, lower expense for a fixed funding
requirement, or a relationship-value objective that includes services beyond
funding? Each question can be legitimate. Only the second is answered by the
funding-cost construction developed here. The construction and all numerical
values in this section are authored, exact illustrations. They supply a
reference for testing a proposed translator, not observed bank economics
\citep{exec17_l2_monetary_exact}.

Let $B_d(a)$ be end-of-day eligible balance under offer $a$ on day $d$, where
$a=0$ is the stated existing offer and $a=1$ is the proposed offer. Fix a
horizon of $H$ days and a day-count denominator $D$. Define
\begin{equation}
 \overline B(a)=\frac1H\sum_{d=1}^H B_d(a),\qquad
 S(a)=\sum_{d=1}^H B_d(a)=H\overline B(a).
 \label{exec17:l2-stock}
\end{equation}
The units are dollars and dollar-days respectively. With a constant annual
valuation rate $q$, the monetary amount assigned to the balance-time is
$qS(a)/D$ dollars. The rate, denominator, daily measurement convention and
horizon are part of this definition. For example, substituting a last-day
balance for an average balance changes the outcome even if the two amounts
happen to agree in a stable illustrative account. A genuine proposal needs
the product's actual accrual convention; the example fixes $D=365$ without
claiming that this convention applies to every deposit contract.

Suppose, for the moment, that treasury has the same exogenous daily funding
requirement $F_d$ under either offer and can cover the residual $F_d-B_d(a)$
at a supplied rate $q_d$. Let $r_d(a)$ be the deposit rate, and let $K(a)$ be
the monetary expense outside interest: for example a stipulated communication
charge. Ignoring discounting in this short-horizon construction, define
\begin{align}
 N(a)&=\frac1D\sum_{d=1}^H
 \{r_d(a)B_d(a)+q_d[F_d-B_d(a)]\}+K(a),
 \label{exec17:l2-cost}\\
 V(a)&=\frac1D\sum_{d=1}^H[q_d-r_d(a)]B_d(a)-K(a).
 \label{exec17:l2-value}
\end{align}
Here $N$ is net funding cost and $V$ is the saving relative to funding the
entire requirement at $q_d$. Their sum is the common reference cost
$D^{-1}\sum_d q_dF_d$. Therefore the offer's incremental value is exactly
$\Delta V=\mathbb E[V(1)-V(0)]=-\mathbb E[N(1)-N(0)]$.
Maximizing $V$ and minimizing $N$ agree. Neither convention licenses
maximizing the increase in funding cost. A separate expense must be netted
as $K(1)-K(0)$, not automatically as $K(1)$: the comparator itself may incur
expense, and cancelling a campaign may save it.

This identity depends on the reference funding requirement and valuation
path being common. If the offer changes lending volumes, the institution's
marginal funding rate, or the capacity to raise substitute funds, those
objects acquire potential-outcome indices and the common-reference term no
longer cancels. An assistant cannot preserve the simple identity while
quietly adding those effects to its narrative. It must expand the target,
state the additional intervention and equilibrium assumptions, or present
the fixed-requirement result as a narrower answer. Likewise, relationship
benefits should be represented by a separately specified monetary outcome,
not attached to a balance effect through an unexplained positive adjective.

\subsection{Repricing existing balances can reverse the conclusion}

Under constant rates, write $r_1=r_0+\delta r$ and
$\Delta S=S(1)-S(0)$. Subtracting the two value definitions gives the
pathwise identity
\begin{equation}
 V(1)-V(0)=\frac{(q-r_0)\Delta S-\delta r S(1)}D
             -[K(1)-K(0)].
 \label{exec17:l2-decomposition}
\end{equation}
The first term values additional balance-time at the old spread. The second
charges the higher deposit rate to all balances receiving it, including
balances that would have remained without the offer. It is not restricted
to newly attracted deposits. An equivalent decomposition uses the new
spread on $\Delta S$ and the rate increase on $S(0)$. Both produce the same
cash difference; adding both repricing terms would count part of the cost
twice. These identities hold for each pair of potential balance paths, so
taking expectations requires no independence between the two paths.

Consider an exact thirty-day example. The comparator holds USD100,000 every
day at a 2\% annual rate; the offered rate is 2.25\%, and its balance is
USD110,000 every day. The stipulated substitute funding rate is 5\%.
Operating expenses are USD4 under the comparator and USD8 under the offer.
Thus $\Delta S=300{,}000$ dollar-days and the added average balance is
USD10,000. The additional old-spread value is $1800/73$ dollars and the
repricing charge is $1650/73$ dollars. Their difference is only $150/73$
dollars. Subtracting the incremental USD4 expense gives
\begin{equation}
 \Delta V=\frac{1800-1650}{73}-4=-\frac{142}{73}
          \simeq -\mathrm{USD}\,1.95.
 \label{exec17:l2-numeric}
\end{equation}
The offer attracts balances while increasing net funding cost. The reference
also constructs the complete daily ledger for a stipulated USD150,000 funding
requirement and reproduces the opposite-signed cost difference by summation.
That second calculation checks the accounting identity instead of merely
repeating the printed formula. Treating the USD10,000 stock increase as a
USD10,000 benefit, forgetting existing-balance repricing, or subtracting
USD8 instead of the incremental USD4 each answers a different question.

The all-eligible mean effect remains a population target. A deterministic
balance path here makes the arithmetic transparent; a causal analysis must
identify or bound the distribution of paths under both offers. Identifying
the mean balance contrast alone is sufficient for this linear, fixed-rate
valuation only when the other required expectations and costs are supplied.
It is not a claim that a regression coefficient identifies the complete
financial response of a deposit portfolio.

\subsection{Marginal funding cost and an explicit valuation boundary}

Another reasonable business question is the annualized incremental expense
per unit of extra balance-time. For a positive expected balance-time change,
under the same constant-rate convention define
\begin{equation}
 r_{\mathrm{inc}}=
 \frac{r_1\mathbb E[S(1)]-r_0\mathbb E[S(0)]
       +D\mathbb E[K(1)-K(0)]}
      {\mathbb E[S(1)-S(0)]}.
 \label{exec17:l2-marginal}
\end{equation}
This is a ratio of portfolio expectations, not the mean of customer-level
ratios. The latter can be undefined for accounts with no balance response,
give large influence to tiny denominators, and change the aggregation
question. The incremental cost rate also differs from the offered deposit
rate. In the example it is $1571/30000$, or about 5.237\%, despite the
2.25\% offered rate. It includes paying more on existing balances and
spreading the incremental operating expense over the new balance-time.
With positive expected $\Delta S$, $\Delta V$ is positive precisely when
$q>r_{\mathrm{inc}}$. If that denominator is zero or negative, the displayed
ratio is unavailable or loses this comparison interpretation; the original
monetary difference remains the appropriate object to inspect.

This comparison is useful even when a bank has not yet supplied a valuation
rate. Instead of inventing a price, the proposal can report the break-even
rate or a sensitivity interval. If the same example permits $q$ between
4\% and 6\%, its value ranges from $-742/73$ to $458/73$ dollars, roughly
minus USD10.16 to plus USD6.27. There is no uniformly preferred offer over
that stipulated set. An owner who accepts the funding model can resolve
the economic choice by supplying the relevant rate and cost convention.
This is a missing business input, distinct from an unidentified balance
response. Uncertainty about both should be kept visible rather than folded
into one unexplained error bar.

For a set of admissible causal laws $\mathcal P$ and valuation specifications
$\mathcal Q$, the corresponding proposal can report
\begin{equation}
 \mathcal V=\{\mathbb E_P[V_q(1)-V_q(0)]:(P,q)\in\mathcal A\},
 \qquad \mathcal A\subseteq\mathcal P\times\mathcal Q.
 \label{exec17:l2-joint}
\end{equation}
The joint admissible set $\mathcal A$ matters. If liquidity conditions affect
both the response and the marginal rate, independently combining marginal
endpoints may create impossible scenarios or omit admissible dependence.
The interval above deliberately varies only a deterministic rate while
holding the response law fixed. It is an economic sensitivity calculation,
not a confidence interval, a Bayesian posterior, or a general sharp bound.
The formulation should state which of those objects the user is requesting
before selecting a downstream numerical tool.

\subsection{Timing and nonlinear valuation require more than a mean}

Balance-time suffices only for a constant linear valuation. With a supplied
time-varying path $q_d$, the gross monetary effect is
$D^{-1}\sum_d q_d\mathbb E[B_d(1)-B_d(0)]$. As an exact two-day example,
compare additional balances of $(200,0)$ with $(0,200)$ dollars. Both have
average additional balance USD100 and exposure 200 dollar-days. At rates
$(0.10,0.01)$ their gross values are $20/365$ and $2/365$ dollars. A translator
that passes only the mean balance to the valuation engine loses the
information needed to distinguish these cases. Discount factors, when
required, must similarly attach to the appropriate cash-flow dates; they
cannot be appended after aggregation if the timing information was dropped.

A nonlinear valuation needs still more information. Suppose a hypothetical
gross value rule is $m(B)=0.01\min(B,100)$ dollars and a comparator has zero
gross value. A balance fixed at USD100 has expected gross value USD1.
A balance of zero or USD200 with equal probabilities has the same mean
balance but expected gross value USD0.50. With incremental expense USD0.75,
their net values are USD0.25 and minus USD0.25. The mean balance response
does not identify which decision is favorable. This example is deliberately
stylized; its purpose is to distinguish $\mathbb E[m(B)]$ from
$m(\mathbb E[B])$, not to suggest a bank's valuation schedule.

Where costs depend on utilization, default, or other post-offer quantities,
those quantities belong in the monetary potential outcome. Their joint
distribution can matter. Using the realized balance to calculate interest
is legitimate outcome construction; using that same post-offer balance as
a pre-offer targeting feature is a clock violation. Similarly, conditioning
on it to estimate an ordinary total effect generally changes the question.
The declaration must identify the role a variable plays, rather than
labelling a field universally permitted or forbidden. The target-trial and
time-varying treatment distinctions already developed from
\citet{hernan2020} provide the inherited framework; the cash-flow mappings
here are project-specific definitions.

\subsection{Consolidation and feasibility are part of the intervention}

The business unit determines what counts as additional funding. If an offer
moves USD10,000 from one account to another account at the same bank, the
receiving account's effect is positive while the consolidated balance change
is zero. Repricing the receiving account can nevertheless increase expense.
For mutually exclusive account balances under a fixed consolidation map,
the bank-level outcome sums the account paths before applying the bank-level
valuation. Shared accounts must not be counted once for each linked owner.
When the intervention changes ownership, currency composition or legal
availability, a fixed consolidation map itself becomes an assumption to
reconsider. An identifier matching exercise cannot choose the economic
boundary on the owner's behalf.

This also changes treatment definition. A rate offered to one account may
alter another account's balance or a household's other offers. The potential
outcome may then depend on an allocation vector, not just an account's own
binary action. A feasible policy must respect the shared budget and the
permitted unit of allocation. Account-level exchangeability and a
no-interference assumption are substantive restrictions on that target,
not conveniences a language model can infer from a rectangular dataset.
If these restrictions are not defensible, the proposal can retain a cluster
or portfolio question with its additional requirements.

A budget needs an equally clear comparator. An absolute expense constraint
$\mathbb E[K(\pi)]\leq B$ differs from an incremental constraint
$\mathbb E[K(\pi)-K(\pi_0)]\leq B$. A capacity constraint on the number of
contacts is a third object and may not be expressible as either monetary
constraint. The policy $\pi$ can use only the declared predecision history;
the comparator $\pi_0$ must satisfy its own declared feasibility conditions.
The proposed assistant should present these alternatives in the business
brief and carry the accepted choice into the formal query. Optimization
before that choice can produce a precise answer to an unstated problem.

\section{Measurement and selective notes can change the target}
\label{exec17:l2-measurement}

The cash-flow argument generalizes to language-derived outcomes. A rule
$g$ that converts a servicing note into a distress indicator defines an
observed measurement, not automatically the underlying distress variable.
The inspected discussions by Keith and colleagues and Egami and colleagues
separate measurement from causal identification and distinguish a learned
text representation from the scientific target
\citep{exec13_census_text_review_arxiv,exec13_census_texts}. The proposed
translator must say whether the owner wants the effect on recorded distress,
on a fixed coding rule's output, or on latent distress. Those outcomes can
differ even when the extraction system executes exactly as specified.

Let a binary underlying outcome have intervention-specific mean
$\mu_a=\Pr(Y(a)=1)$. For a fixed binary measurement rule, define its
sensitivity $s_a=\Pr(Y^*=1\mid Y=1,\operatorname{do}(A=a))$ and specificity
$c_a=\Pr(Y^*=0\mid Y=0,\operatorname{do}(A=a))$. By the law of total
probability, when these conditional rates are defined,
\begin{equation}
 \mu_a^*=s_a\mu_a+(1-c_a)(1-\mu_a)
          =(s_a+c_a-1)\mu_a+1-c_a.
 \label{exec17:l2-error}
\end{equation}
This is an exact project derivation, not an imported guarantee for an LLM.
If sensitivity and specificity are common across interventions, the measured
contrast equals $(s+c-1)(\mu_1-\mu_0)$. Recovering the underlying contrast
by division requires those error rates to be known, transportable and
$s+c-1\ne0$. If the sum is one, the measurement contains no marginal
information about the binary outcome through this channel. If it is less
than one, its direction is reversed. Assuming ordinary attenuation without
checking the channel is therefore unjustified.

For example, set true outcome risks to $2/5$ and $1/2$, a true effect of
$1/10$. With common sensitivity $4/5$ and specificity $9/10$, the measured
effect is $7/100$. Now retain those control error rates but use treated
sensitivity $1/2$ and specificity $19/20$. The measured effect becomes
$-21/200$, reversing the sign. All probabilities remain valid. The change
could represent different document content or extraction errors after the
offer; the example makes no claim that this occurs in a bank. Its implication
is that pooled classification accuracy cannot determine the causal error.
A validation design needs outcome truth and enough treatment-specific
information to support the proposed correction, including the population
and document mechanism to which the rates are transported.

Missing notes introduce a separate problem even when every available note
is coded perfectly. Let $A$ and predecision distress $U$ be independent
Bernoulli variables with probability $1/2$, and define $Y(a)=U$ for either
action. The intervention has zero effect. Let $R$ indicate a note available
for analysis, with probability $3/4$ when $A=U$ and $1/4$ otherwise.
Every action--distress combination has positive note availability, so this
example does not depend on a cell with zero support. Exact enumeration gives
\begin{equation}
 \Pr(Y=1\mid A=1,R=1)=\frac34,\qquad
 \Pr(Y=1\mid A=0,R=1)=\frac14.
 \label{exec17:l2-selection}
\end{equation}
The observed-note contrast is $1/2$ despite random assignment and a zero
effect in the full population. Selecting notes conditions on a common effect
of action and distress. A correct event timestamp, perfect extraction and
randomized treatment do not remove that selection mechanism.

An oracle that knows $U$ and the note probabilities can reweight each
selected cell by the inverse of its selection probability and recover
mean distress $1/2$ in each arm. This verifies the construction; it is not
a proposed correction when distress is precisely the unobserved quantity
the notes were meant to measure. The useful handoff identifies what is
missing: a justified observation model, independent validation measurements,
or a narrower target that the owner actually accepts. It must not quietly
replace the eligible population by customers with convenient documentation.

These examples distinguish three repairs. Correcting a unit or comparator
is a specification repair. Estimating a measurement channel needs suitable
validation information. Correcting selective availability needs a defensible
observation law and support. A language model can locate definitions,
propose alternatives and flag contradictions in each case, but those actions
do not provide the missing scientific information. Multiple defensible
diagrams and their different identifying consequences should survive the
formulation stage, consistent with the inspected argument for explicit
assumption-based modelling in \citet{exec13_census_credible}.

\subsection{A proposal evaluation that isolates the translation error}

A useful public and generated evaluation can freeze a set of business
briefs that share language while differing in one decision-relevant fact.
For funding, pair a balance target with a net-cost target; then vary whether
the comparator incurs expense, whether repricing covers old balances, and
whether valuation is time-varying or nonlinear. For measurement, vary
common versus differential error and complete versus selective notes.
The exact mechanisms above supply the expected mathematical answer, while
public dictionaries can supply realistic field names and documented units.
Public text alone supplies neither assignment truth nor a bank's marginal
funding rate. An evaluation should preserve that distinction explicitly.

The reference response should contain the target, the necessary outcome
components, the sign convention, the information that is absent, and the
reason an alternative interpretation would change the answer. A numerical
score is appropriate only when the brief supplies enough information to
determine it. For an unspecified $q$, reporting the break-even rate and the
unresolved valuation choice can be correct; inventing a rate to produce a
confident investment recommendation is an error. For selective notes with
unknown selection probabilities, reporting the selected contrast as the
full-population effect is an error even when the subtraction is exact.

The comparison should use the same briefs, dictionaries, permissible source
material and answer schema for a structured analyst baseline and an assisted
workflow. Separate the underlying mechanisms used for development from those
used for final evaluation, and keep ambiguous, unsupported and failed requests
in the denominator. This section provides exact authored reference families,
not a completed provider comparison. A useful prototype result would identify
which translation failures become less frequent and at what review cost.
That result can support a later decision to purchase authentic banking
validation; it cannot supply such validation in advance.
\label{exec17:l2-monetary-end}
